% from ega4-2-fr.tex
\oldpage[IV-2]{139}

en outre, comme $N$ est un $A$-module plat, l'hypothèse que $x$ est
$M$-régulier entraîne que $x$ est aussi $(M\otimes_AN)$-régulier
$(0_{\mathrm I}, 6.1.1)$. On a par suite (0, 16.4.6, (ii) et 16.4.8)
\[
\operatorname{prof}_{A'}(M')=\operatorname{prof}_A(M)-1,
\qquad
\operatorname{prof}_{B'}(M'\otimes_{A'}N')
=\operatorname{prof}_B(M\otimes_AN)-1
\]
et
\[
\operatorname{prof}_{B'\otimes_{A'}k}(N'\otimes_{A'}k)
=\operatorname{prof}_{B\otimes_Ak}(N\otimes_Ak).
\]

L'égalité (6.3.1.1) est donc conséquence de la même relation pour $A'$,
$B'$, $M'$ et $N'$; mais comme $N$ est un $A$-module plat,
$N'=N\otimes_AA'$ est un $A'$-module plat $(0_{\mathrm I}, 6.2.1)$; on peut
par suite appliquer l'hypothèse de récurrence, ce qui démontre (6.3.1.1) dans
ce cas.

$b)$ Supposons $\operatorname{prof}_{B\otimes_Ak}(N\otimes_Ak)>0$.
Considérons un élément $y\in\mathfrak m$ qui soit
$(N\otimes_Ak)$-régulier; il résulte de $(0_{\mathrm{III}}, 10.2.4)$ que $y$
est alors $N$-régulier et que
\[
N'=N/yN
\]
est un $A$-module plat, puisque $N$ est supposé être un $A$-module plat.
Appliquant $(0_{\mathrm I}, 6.1.2)$ à la suite exacte de $A$-modules
\[
0\to N\xrightarrow{y}N\to N'\to0
\]
on en conclut des isomorphismes
\[
(M\otimes_AN)/y(M\otimes_AN)\xrightarrow{\sim}M\otimes_AN'
\quad\text{et}\quad
N'\otimes_Ak\xrightarrow{\sim}(N\otimes_Ak)/y(N\otimes_Ak)
\]
et en outre que $y$ est $(M\otimes_AN)$-régulier. On a par suite
\[
\operatorname{prof}_{B\otimes_Ak}(N'\otimes_Ak)
=\operatorname{prof}_{B\otimes_Ak}(N\otimes_Ak)-1
\]
et
\[
\operatorname{prof}_B(M\otimes_AN')
=\operatorname{prof}_B(M\otimes_AN)-1;
\]
l'hypothèse de récurrence montre que la relation (6.3.1.1) est valable pour
$A$, $B$, $M$ et $N'$, et d'après ce qui précède, on en déduit (6.3.1.1)
pour $A$, $B$, $M$ et $N$.

\begin{corollary}[6.3.2]
\label{IV.2.s6.3.2-fr}
Sous les hypothèses de (6.3.1), supposons en outre $M\neq0$ et $N\neq0$;
alors on a
\begin{equation}
\operatorname{coprof}_B(M\otimes_AN)=\operatorname{coprof}_A(M)
+\operatorname{coprof}_{B\otimes_Ak}(N\otimes_Ak).
\tag{6.3.2.1}\label{IV.2.s6.3.2.1-fr}
\end{equation}
\end{corollary}

Cela résulte aussitôt de (6.3.1.1) et (6.1.2) et de la définition de la
coprofondeur (0, 16.4.9).

\begin{corollary}[6.3.3]
\label{IV.2.s6.3.3-fr}
Sous les hypothèses de (6.3.1), supposons en outre $M\neq0$ et $N\neq0$; pour
que $M\otimes_AN$ soit un $B$-module de Cohen-Macaulay, il faut et il suffit
que $M$ soit un $A$-module de Cohen-Macaulay et que $N\otimes_Ak$ soit un
$(B\otimes_Ak)$-module de Cohen-Macaulay.
\end{corollary}

Cela résulte du corollaire (6.3.2) et de la définition des modules de
Cohen-Macaulay par le fait que leur coprofondeur est nulle (0, 16.5.1), tenant
compte de ce que la condition $N\neq0$ est équivalente à $N\otimes_Ak\neq0$ en
vertu du lemme de Nakayama.

\begin{corollary}[6.3.4]
\label{IV.2.s6.3.4-fr}
Sous les hypothèses de (6.3.1), supposons en outre que $N\otimes_Ak$ soit un
$B$-module de longueur finie; alors on a
\begin{equation}
\operatorname{prof}_B(M\otimes_AN)=\operatorname{prof}_A(M).
\tag{6.3.4.1}\label{IV.2.s6.3.4.1-fr}
\end{equation}

